\section*{Capítulo \ref{ch:b1:structures} --- Estruturas Algébricas}

\begin{solution}{exo:b1:structures:1}
\emph{Estabilidade:} $x * y = 1 \iff x + y - xy = 1 \iff (1-x)(1-y) = 0$,
impossível para $x, y \neq 1$. Com efeito, a identidade-chave é
\[
1 - x * y = (1 - x)(1 - y):
\]
a \omterm{def:b1:logic:map}{aplicação} $\varphi(x) = 1 - x$ leva $(E, *)$ em $(\R^*, \times)$ com
$\varphi(x * y) = \varphi(x)\varphi(y)$ --- um \omterm{def:b1:structures:morphism}{morfismo} \omterm{def:b1:logic:inj}{bijetivo}. Todos os
axiomas se transportam agora: a associatividade e a comutatividade
decorrem das de $\times$; o neutro é $\varphi^{-1}(1) = 0$ (verificação: $x *
0 = x$); o inverso de $x$ é $\varphi^{-1}\bigl((1-x)^{-1}\bigr) =
1 - \frac{1}{1-x} = \frac{x}{x - 1}$ (que é $\neq 1$). Logo, $(E, *)$
é um \omterm{def:b1:structures:group}{grupo abeliano}.
\end{solution}

\begin{solution}{exo:b1:structures:2}
\begin{enumerate}
  \item Sim: o produto de positivos é positivo, o neutro é $1$, o inverso
        é $\frac 1x$, e a associatividade é herdada de $\R^*$.
  \item Não: não é estável ($1 + 1 = 2 \notin \{-1,0,1\}$).
  \item Sim: o exemplo padrão.
  \item Não: não é estável (ímpar $+$ ímpar $=$ par) e não há neutro ($0$
        é par).
\end{enumerate}
\end{solution}

\begin{solution}{exo:b1:structures:3}
Escreva $\mathrm{id}$, as transposições $\tau_{12}, \tau_{13},
\tau_{23}$ (que trocam os dois pontos nomeados) e os ciclos $c =
(1\,2\,3)$ (isto é, $1 \mapsto 2 \mapsto 3 \mapsto 1$) e $c^2 =
(1\,3\,2)$. A tabela de $\sigma\rho$ (linha $\sigma$, coluna $\rho$,
aplicando $\rho$ primeiro):
\begin{center}
\renewcommand{\arraystretch}{1.2}
\begin{tabular}{c|cccccc}
$\sigma\backslash\rho$ & $\mathrm{id}$ & $c$ & $c^2$ & $\tau_{12}$ &
$\tau_{13}$ & $\tau_{23}$ \\
\hline
$\mathrm{id}$ & $\mathrm{id}$ & $c$ & $c^2$ & $\tau_{12}$ &
$\tau_{13}$ & $\tau_{23}$ \\
$c$ & $c$ & $c^2$ & $\mathrm{id}$ & $\tau_{13}$ & $\tau_{23}$ &
$\tau_{12}$ \\
$c^2$ & $c^2$ & $\mathrm{id}$ & $c$ & $\tau_{23}$ & $\tau_{12}$ &
$\tau_{13}$ \\
$\tau_{12}$ & $\tau_{12}$ & $\tau_{23}$ & $\tau_{13}$ & $\mathrm{id}$
& $c^2$ & $c$ \\
$\tau_{13}$ & $\tau_{13}$ & $\tau_{12}$ & $\tau_{23}$ & $c$ &
$\mathrm{id}$ & $c^2$ \\
$\tau_{23}$ & $\tau_{23}$ & $\tau_{13}$ & $\tau_{12}$ & $c^2$ & $c$ &
$\mathrm{id}$ \\
\end{tabular}
\end{center}
Par que não comuta: $\tau_{12}\tau_{13} = c^2$, ao passo que
$\tau_{13}\tau_{12} = c$. (Para conferir uma entrada:
$\tau_{12}\tau_{13}$ envia $1 \xmapsto{\tau_{13}} 3
\xmapsto{\tau_{12}} 3$, $3 \mapsto 1 \mapsto 2$, $2 \mapsto 2 \mapsto
1$: isto é, $1 \mapsto 3 \mapsto 2 \mapsto 1$, o ciclo $c^2 =
(1\,3\,2)$.)
\end{solution}

\begin{solution}{exo:b1:structures:4}
$H$: $1 \in H$; para $z, w \in H$, $\abs{zw^{-1}} = \abs z / \abs w =
1$: o critério se aplica. $\R_+^*$: idem, com $\abs{xy^{-1}}$ substituído
pela positividade. União: $\iu \in H$ e $2 \in \R_+^*$, mas $2\iu$ tem
módulo $2 \neq 1$ e não é um real positivo: $2\iu \notin H \cup
\R_+^*$, de modo que a união não é estável --- não é \omterm{def:b1:structures:subgroup}{subgrupo} (como
previsto pelo \cref{exo:b1:structures:6}, já que nenhum dos \omterm{def:b1:structures:subgroup}{subgrupos} contém o outro).
\end{solution}

\begin{solution}{exo:b1:structures:5}
\omterm{def:b1:structures:morphism}{Morfismo}: $\eu^{\iu(\theta + \varphi)} =
\eu^{\iu\theta}\eu^{\iu\varphi}$ (\cref{thm:b1:complex:funceq}).
\omterm{def:b1:structures:morphism}{Núcleo}: $\eu^{\iu\theta} = 1 \iff \theta \in 2\pi\Z$, de modo que $\ker f =
2\pi\Z \neq \{0\}$: não \omterm{def:b1:logic:inj}{injetivo}. Imagem: todo número complexo de módulo $1$ é
$\eu^{\iu\theta}$ para algum $\theta$ (forma polar), de modo que
$\operatorname{im} f = \mathbb{U}$, o círculo unitário. A restrição
de $f$ a um intervalo semiaberto de comprimento $2\pi$, como
$\intco{0}{2\pi}$ ou $\intoc{-\pi}{\pi}$, é \omterm{def:b1:logic:inj}{injetiva} (dois ângulos
com a mesma imagem diferem por um múltiplo de $2\pi$, e só um
representante de cada classe cabe no intervalo); nenhum intervalo de
comprimento maior funciona, pois ele contém dois pontos à distância
$2\pi$.
\end{solution}

\begin{solution}{exo:b1:structures:6}
Interseção: $e \in H \cap K$, e $x, y \in H \cap K$ dá
$xy^{-1}$ tanto em $H$ quanto em $K$. União: se $H \subseteq K$, a união é
$K$, um \omterm{def:b1:structures:subgroup}{subgrupo} (e simetricamente). Reciprocamente, suponha que nenhuma
das inclusões valha: tome $h \in H \setminus K$ e $k \in K \setminus
H$, e suponha que $H \cup K$ fosse um \omterm{def:b1:structures:subgroup}{subgrupo}; então $hk \in H \cup K$.
Se $hk \in H$, então $k = h^{-1}(hk) \in H$: contradição. Se $hk \in
K$, então $h = (hk)k^{-1} \in K$: contradição. Logo, $H \cup K$ não é
\omterm{def:b1:structures:subgroup}{subgrupo}.
\end{solution}

\begin{solution}{exo:b1:structures:7}
Note primeiro que $x^2 = e$ significa $x^{-1} = x$ para todo $x$. Então, para
$x, y \in G$:
\[
xy = (xy)^{-1} = y^{-1} x^{-1} = yx ,
\]
usando a \cref{prop:b1:structures:rules}~(2). Logo, $G$ é \omterm{def:b1:structures:group}{abeliano}.
\end{solution}

\begin{solution}{exo:b1:structures:8}
Unidades de $\Z/18\Z$: as classes primas com $18 = 2 \times 3^2$:
$\overline 1, \overline 5, \overline 7, \overline{11}, \overline{13},
\overline{17}$. Inversa de $\overline 5$: $5 \times 11 = 55 = 3\times
18 + 1$, de modo que $\overline 5^{-1} = \overline{11}$.

$\overline 5 x = \overline 7$: multiplique por $\overline{11}$: $x =
\overline{77} = \overline 5$ (pois $77 = 4\times 18 + 5$). Solução
única.

$\overline 6 x = \overline 3$: a equação $6x \equiv 3 \pmod{18}$
significa que $18 \mid 6x - 3$. Mas $6x - 3 = 3(2x - 1)$ é ímpar, ao passo que $18$
é par: um número par não pode dividir um ímpar. Sem solução.

$\overline 6 x = \overline{12}$: $6x \equiv 12 \pmod{18} \iff x
\equiv 2 \pmod 3$: soluções $x \in \{\overline 2, \overline 5,
\overline 8, \overline{11}, \overline{14}, \overline{17}\}$ --- seis
delas.
\end{solution}

\begin{solution}{exo:b1:structures:9}
$\Z[\sqrt 2]$ contém $0$ e $1$, e é estável por subtração
e por produto:
\[
(a + b\sqrt 2)(c + d\sqrt 2) = (ac + 2bd) + (ad + bc)\sqrt 2 ,
\]
de modo que é um subanel de $\R$ (comutatividade, associatividade e
distributividade são herdadas). Unidade: $(1 + \sqrt 2)(-1 + \sqrt 2) =
2 - 1 = 1$, de modo que $1 + \sqrt 2$ é invertível, com inversa $\sqrt 2 - 1
\in \Z[\sqrt 2]$. As suas potências $(1 + \sqrt 2)^n$ são estritamente
crescentes (a base é $> 1$), logo duas a duas distintas, e cada uma é
uma unidade ($\bigl((1+\sqrt2)^n\bigr)^{-1} = (\sqrt 2 - 1)^n$): o \omterm{def:b1:structures:group}{grupo}
das unidades é infinito.
\end{solution}

\begin{solution}{exo:b1:structures:10}
$x + x = (x + x)^2 = x^2 + x^2 + x^2 + x^2 = 4x^2 = 4x$ --- de modo que $2x =
4x$, o que dá $2x = 0$, isto é, $x + x = 0$ (cada elemento é o seu próprio
inverso aditivo). Então
\[
x + y = (x+y)^2 = x^2 + xy + yx + y^2 = x + xy + yx + y ,
\]
de modo que $xy + yx = 0$, isto é, $xy = -yx = yx$ (usando $-z = z$). Portanto, $A$ é
comutativo.

Exemplo: em $\mathcal{P}(E)$, defina $A + B = (A \cup B) \setminus (A
\cap B)$ (diferença simétrica) e $A \times B = A \cap B$. Verifica-se
que: $(\mathcal{P}(E), +)$ é um \omterm{def:b1:structures:group}{grupo abeliano} com neutro
$\emptyset$ e cada \omterm{def:b1:logic:sets}{conjunto} é o seu próprio inverso; $\cap$ é associativa,
comutativa, com neutro $E$; a distributividade $A \cap (B + C) = (A
\cap B) + (A \cap C)$ vale (um elemento está no lado esquerdo se, e somente se,
está em $A$ e em exatamente um dentre $B, C$). E $A \cap A = A$: todo
elemento é idempotente, como exigido.
\end{solution}

\begin{solution}{exo:b1:structures:11}
Escreva a hipótese para $n$ e $n+1$:
\[
(xy)^{n+1} = x^{n+1} y^{n+1}
\quad\text{e}\quad
(xy)^{n+1} = (xy)(xy)^n = xy\,x^n y^n .
\]
Igualando: $x^{n+1} y^{n+1} = x\,y\,x^n\,y^n$; cancele $x$ à esquerda
e $y^n$ à direita: $x^n y = y x^n$. O mesmo cálculo um
grau acima ($n+1$ e $n+2$) dá $x^{n+1} y = y x^{n+1}$. Então
\[
y\,x^{n+1} = x^{n+1} y = x\,(x^n y) = x\,y\,x^n ,
\]
e, cancelando $x^n$ à direita em $y x \cdot x^n = x y \cdot x^n$:
$yx = xy$. Logo, $G$ é \omterm{def:b1:structures:group}{abeliano}.
\end{solution}

\begin{solution}{exo:b1:structures:12}
\begin{enumerate}
  \item Seja $f \colon \Z \to \Z$ aditiva e $a = f(1)$. Por
        indução, $f(k) = ka$ para $k \in \N$, e $f(-k) = -f(k) =
        -ka$: logo, $f$ é a multiplicação por $a$. Reciprocamente,
        toda \omterm{def:b1:logic:map}{aplicação} $k \mapsto ak$ é um \omterm{def:b1:structures:morphism}{morfismo}: os \omterm{def:b1:structures:morphism}{morfismos}
        $(\Z,+) \to (\Z,+)$ são exatamente as multiplicações por um
        inteiro fixado.
  \item Seja $f \colon \Q \to \Z$ um \omterm{def:b1:structures:morphism}{morfismo}, $x \in \Q$ e
        $n \in \N^*$. Então
        \[
        f(x) = f\Bigl(\underbrace{\tfrac xn + \dots +
        \tfrac xn}_{n}\Bigr) = n\,f\Bigl(\frac xn\Bigr) ,
        \]
        de modo que o inteiro $f(x)$ é divisível por todo $n \geq 1$.
        O único inteiro assim é $0$: $f \equiv 0$.
\end{enumerate}
\end{solution}

\begin{solution}{pb:b1:structures:1}
\textbf{1.} Uma \omterm{def:b1:counting:objects}{permutação} é uma bijeção de $\intint1n$, isto é, um
$n$-arranjo de $n$ objetos: há $n!$ delas
(\cref{thm:b1:counting:counts}). Com $\sigma = [2,3,1]$, $\tau =
[1,3,2]$: $\sigma\tau$ envia $1 \mapsto 1 \mapsto 2$, $2 \mapsto 3
\mapsto 1$, $3 \mapsto 2 \mapsto 3$: $\sigma\tau = [2,1,3]$; e
$\tau\sigma$ envia $1 \mapsto 2 \mapsto 3$, $2 \mapsto 3 \mapsto
2$, $3 \mapsto 1 \mapsto 1$: $\tau\sigma = [3,2,1] \neq
\sigma\tau$.

\textbf{2.} Sejam $\gamma, \gamma'$ com suportes disjuntos $S, S'$.
Para $x \in S$: $\gamma'(x) = x$ e $\gamma(x) \in S$, de modo que
$\gamma\gamma'(x) = \gamma(x) = \gamma'\gamma(x)$. Simetricamente
para $x \in S'$; e os dois lados fixam todo $x \notin S \cup S'$. Logo,
$\gamma\gamma' = \gamma'\gamma$.

\textbf{3.} Para $i \in \intint1n$, os valores $i, \sigma(i),
\sigma^2(i), \dots$ vivem num \omterm{def:b1:counting:card}{conjunto finito}, de modo que $\sigma^a(i) =
\sigma^b(i)$ para certos $a < b$; a injetividade dá $\sigma^{b-a}(i)
= i$: a sequência volta a $i$. Chame de \emph{órbita} de $i$ o
\omterm{def:b1:logic:sets}{conjunto} $\{i, \sigma(i), \dots, \sigma^{k-1}(i)\}$, com $k \geq 1$
mínimo tal que $\sigma^k(i) = i$. Duas órbitas que se encontram num
ponto coincidem (as duas são as imagens de $\sigma$ para a frente desse
ponto), de modo que as órbitas formam uma \omterm{thm:b1:logic:partition}{partição} de $\intint1n$; $\sigma$
age em cada órbita de tamanho $k \geq 2$ como o $k$-ciclo $(i\ \sigma(i)\ \cdots\
\sigma^{k-1}(i))$ e fixa os pontos isolados. O produto desses
ciclos disjuntos coincide com $\sigma$ em toda parte. Unicidade: em
qualquer decomposição em ciclos disjuntos, o ciclo que passa por $i$ tem
de ser $(i\ \sigma(i)\ \cdots)$ --- os ciclos ficam forçados a serem as
órbitas com a ação induzida.

\textbf{4.} Seguindo as órbitas: $1 \to 4 \to 2 \to 1$, $3 \to 5
\to 3$, $6 \to 7 \to 8 \to 6$:
\[
\sigma = (1\ 4\ 2)(3\ 5)(6\ 7\ 8) .
\]
Se $\sigma = \gamma_1\cdots\gamma_r$ com ciclos disjuntos de
comprimentos $k_1, \dots, k_r$, a comutação (questão 2) dá
$\sigma^m = \gamma_1^m\cdots\gamma_r^m$ e, como os suportes
são disjuntos, $\sigma^m = \mathrm{id}$ se, e somente se, cada $\gamma_i^m =
\mathrm{id}$, isto é, $k_i \mid m$ para todo $i$ (um $k$-ciclo tem \omterm{def:b1:structures:order}{ordem}
$k$: $\gamma^m$ envia $a_1$ a $a_{1 + (m \bmod k)}$). O menor
$m$ desses é $\operatorname{lcm}(k_1, \dots, k_r)$. Aqui:
$\operatorname{lcm}(3, 2, 3) = 6$.

\textbf{5.} Aplique o lado direito a cada ponto, começando pelo fator
mais à direita. $a_1 \mapsto a_2$ por $(a_1\ a_2)$, e depois todo
fator posterior fixa $a_2$: no total, $a_1 \mapsto a_2$. Para $2 \leq i < k$:
$a_i$ fica intocado até que $(a_1\ a_i)$ o envie a $a_1$, e o
fator imediatamente seguinte $(a_1\ a_{i+1})$ envia $a_1$ a $a_{i+1}$, após
o que nada mais o move: no total, $a_i \mapsto a_{i+1}$. Por fim, $a_k$
é fixado por todos os fatores, salvo o mais à esquerda, que o envia a
$a_1$. Isto é exatamente o ciclo. Como toda \omterm{def:b1:counting:objects}{permutação} é um produto de
ciclos (questão 3), ela é um produto de transposições. Para o $\sigma$
da questão 4:
\[
\sigma = (1\ 2)(1\ 4)\;(3\ 5)\;(6\ 8)(6\ 7),
\]
cinco transposições.

\textbf{6.} Indução em $b - a$. Para $b = a + 1$, a identidade é
trivial ($1 = 2\cdot1 - 1$ fator). Para $b > a + 1$, verifique
diretamente que $(a\ b) = (a\ \ a{+}1)\,(a{+}1\ \ b)\,(a\ \ a{+}1)$:
o lado direito envia $a \mapsto a{+}1 \mapsto b \mapsto b$, $b
\mapsto b \mapsto a{+}1 \mapsto a$, $a{+}1 \mapsto a \mapsto a
\mapsto a{+}1$, e fixa o resto. Por indução, $(a{+}1\ \ b)$ é
um produto palindrômico de $2(b - a - 1) - 1$ transposições
adjacentes, de modo que $(a\ b)$ é um produto de $2(b - a) - 1$: um número
ímpar.

\textbf{7.} $N(\mathrm{id}) = 0$, $\varepsilon = +1$. Para $(i\ \
i{+}1)$, o único par invertido é $(i, i+1)$: $N = 1$,
$\varepsilon = -1$. Para $[2, 3, 1]$: os pares invertidos são $(1,
3)$ (valores $2 > 1$) e $(2, 3)$ (valores $3 > 1$): $N = 2$,
$\varepsilon = +1$.

\textbf{8.} As listas de valores de $\sigma$ e de $\sigma\tau$ diferem
apenas pela troca das posições $i$ e $i + 1$. Para um par de
posições que não envolva $i, i+1$, nada muda. Para $k < i$,
os dois pares $(k, i)$ e $(k, i+1)$ trocam os seus estados de
inversão (os mesmos dois valores são comparados com $\sigma(k)$, na
outra \omterm{def:b1:structures:order}{ordem} de posições): a sua contribuição total fica
inalterada; do mesmo modo para $k > i + 1$. O único par restante $(i,
i+1)$ inverte o seu estado. Portanto, $N(\sigma\tau) = N(\sigma) \pm 1$.

\textbf{9.} Seja $\tau = (a\ b)$ uma transposição qualquer: pela
questão 6 ela é um produto de um número ímpar de transposições
adjacentes, de modo que multiplicar à direita por $\tau$ altera $N$ de
um total ímpar (questão 8, aplicada repetidamente): $\varepsilon(\sigma
\tau) = -\varepsilon(\sigma)$. Agora, se $\sigma = \tau_1\cdots
\tau_p$ (transposições), construa-a a partir da identidade por $p$
multiplicações à direita: $\varepsilon(\sigma) =
(-1)^p\varepsilon(\mathrm{id}) = (-1)^p$. Como
$\varepsilon(\sigma)$ é definido por inversões --- independentemente
de qualquer fatoração --- a paridade de $p$ é um invariante de
$\sigma$. \omterm{def:b1:structures:morphism}{Morfismo}: escrevendo $\sigma$ com $p$ e $\sigma'$ com
$q$ transposições, $\sigma\sigma'$ usa $p + q$ delas:
$\varepsilon(\sigma\sigma') = (-1)^{p+q} =
\varepsilon(\sigma)\varepsilon(\sigma')$.

\textbf{10.} Um $k$-ciclo é um produto de $k - 1$ transposições
(questão 5): $\varepsilon = (-1)^{k-1}$. Para um $\sigma$ geral
com órbitas de tamanhos $k_1, \dots, k_r$ ($k_i \geq 2$) mais $f$
pontos fixos, $c(\sigma) = r + f$ e $n = k_1 + \dots + k_r + f$,
de modo que
\[
\varepsilon(\sigma) = \prod_{i=1}^r (-1)^{k_i - 1}
= (-1)^{\sum_i k_i - r} = (-1)^{n - f - r} = (-1)^{n -
c(\sigma)} .
\]

\textbf{11.} $\mathfrak A_n = \ker\varepsilon$ é um \omterm{def:b1:structures:subgroup}{subgrupo}, por ser o
\omterm{def:b1:structures:morphism}{núcleo} de um \omterm{def:b1:structures:morphism}{morfismo} (\cref{def:b1:structures:morphism}). Fixe
uma transposição $\tau_0$ (existe para $n \geq 2$). A \omterm{def:b1:logic:map}{aplicação} $\sigma
\mapsto \sigma\tau_0$ é uma bijeção de $\mathfrak S_n$ (a sua própria
inversa) que troca $\mathfrak A_n$ pelo \omterm{def:b1:logic:sets}{conjunto} das
permutações ímpares (questão 9). Os dois \omterm{def:b1:logic:sets}{conjuntos} formam uma \omterm{thm:b1:logic:partition}{partição}
de $\mathfrak S_n$ e têm o mesmo tamanho: $\abs{\mathfrak A_n} = \frac{n!}2$.

\textbf{12.} \emph{Inversões} de $[4, 1, 5, 2, 3, 7, 8, 6]$: a partir do
valor $4$: sobre $1, 2, 3$: três; a partir de $5$: sobre $2, 3$: duas;
a partir de $7$: sobre $6$: uma; a partir de $8$: sobre $6$: uma. $N = 7$,
$\varepsilon = -1$. \emph{Tipo de ciclo}: $c = 3$ órbitas, $n = 8$:
$\varepsilon = (-1)^{8-3} = -1$. \emph{Contagem de transposições}: cinco
transposições na questão 5: $(-1)^5 = -1$. As três concordam.

\textbf{13.} $(a\ b)(a\ c)$ (o mais à direita primeiro): $a \mapsto c
\mapsto c$; $c \mapsto a \mapsto b$; $b \mapsto b \mapsto a$:
o $3$-ciclo $(a\ c\ b)$. E $(a\ c\ b)(a\ c\ d)$: $a \mapsto c
\mapsto b$; $b \mapsto b \mapsto a$; $c \mapsto d \mapsto d$; $d
\mapsto a \mapsto c$: isto é, $(a\ b)(c\ d)$, como afirmado.

\textbf{14.} Seja $\sigma \in \mathfrak A_n$: pela questão 9,
$\sigma = \tau_1\cdots\tau_{2m}$ com um número par de
transposições. Agrupe-as em pares consecutivos
$\tau_{2i-1}\tau_{2i}$: se as duas são iguais, o par é a
identidade e desaparece; se elas partilham exatamente um ponto, a
primeira identidade da questão 13 escreve o par como um $3$-ciclo;
se são disjuntas, a segunda identidade o escreve como dois
$3$-ciclos. Portanto, $\sigma$ é um produto de $3$-ciclos (ou a
identidade, um produto vazio --- e, para $n \geq 3$, também $(1\ 2\
3)^3$).

\textbf{15.} $(1\ 2)(3\ 4) = (1\ 3\ 2)(1\ 3\ 4)$ (questão 13
com $a{=}1, b{=}2, c{=}3, d{=}4$). Para o $5$-ciclo: pela
questão 5, $(1\ 2\ 3\ 4\ 5) = (1\ 5)(1\ 4)(1\ 3)(1\ 2)$ e,
emparelhando: $(1\ 5)(1\ 4) = (1\ 4\ 5)$, $(1\ 3)(1\ 2) = (1\ 2\ 3)$:
\[
(1\ 2\ 3\ 4\ 5) = (1\ 4\ 5)(1\ 2\ 3) .
\]
(Verificação em $3$: $(1\ 2\ 3)$ envia $3 \to 1$, e depois $(1\ 4\ 5)$
envia $1 \to 4$: no total, $3 \to 4$, correto.)

\textbf{16.} Aplique os dois lados a um ponto arbitrário. Para $i =
\sigma(a_j)$: o lado esquerdo dá $\sigma\bigl((a_1\ \dots\
a_k)(a_j)\bigr) = \sigma(a_{j+1})$ (índices módulo $k$), que é
o que o lado direito faz com $\sigma(a_j)$. Para $i$ que não seja dessa
forma: $\sigma^{-1}(i)$ está fora do suporte, de modo que o lado esquerdo
fixa $i$, e o lado direito também. Iguais em toda parte.

\textbf{17.} Deslizar a peça da casa $c'$ para a casa vazia
$c$ troca os conteúdos das casas $c$ e $c'$ (a peça $9$, o
vazio, passa para $c'$). Se a peça $\sigma(i)$ estava na casa $i$, a
nova posição é $\sigma' = \sigma \circ (c\ c')$: mesmos conteúdos,
salvo que as casas $c, c'$ passam a ler o antigo conteúdo uma da outra.
Pela questão 9, $\varepsilon(\sigma') = -\varepsilon(\sigma)$.

\textbf{18.} Um movimento leva a casa vazia a uma casa adjacente: a sua
linha ou a sua coluna muda exatamente de $1$, de modo que a distância de
Manhattan $d$ até a casa $9$ muda de $\pm1$, e $(-1)^d$ se inverte. Como
cada movimento inverte tanto $\varepsilon(\sigma)$ quanto $(-1)^{d(\sigma)}$,
o produto deles, $I(\sigma)$, fica inalterado por todo movimento: um
invariante.

\textbf{19.} A posição resolvida tem $\varepsilon = +1$, $d = 0$:
$I = +1$. O alvo (peças $7, 8$ trocadas, vazio em casa) é a
transposição dos conteúdos das casas $7$ e $8$: $\varepsilon
= -1$, $d = 0$: $I = -1$. Como $I$ é invariante e os dois
valores diferem, nenhuma sequência de movimentos os liga.

\textbf{20.} Uma posição com o vazio em casa é uma \omterm{def:b1:counting:objects}{permutação}
das $8$ peças entre as casas $1, \dots, 8$, isto é, um elemento de
$\mathfrak S_8$; ela tem $d = 0$, de modo que $I = \varepsilon(\sigma)$.
Ser alcançável força $I = +1$, isto é, $\sigma \in \mathfrak A_8$; a
recíproca admitida diz que todo $\mathfrak A_8$ é alcançado. Contagem:
$\abs{\mathfrak A_8} = \frac{8!}2 = 20\,160$ (questão 11).

\textbf{21.} Pela questão 20, os arranjos alcançáveis com o vazio em
casa formam exatamente $\mathfrak A_8$ --- em particular, um
\emph{\omterm{def:b1:structures:subgroup}{subgrupo}} de $\mathfrak S_8$: compor dois embaralhamentos solúveis,
ou inverter um deles, continua solúvel, o que está longe de ser
óbvio por puro raciocínio sobre o quebra-cabeça.

\textbf{22.} (a) Um $3$-ciclo de peças com o vazio em casa:
$\varepsilon = +1$ (questão 10), $d = 0$, logo $I = +1$: alcançável
(pela recíproca admitida) --- é possível ciclar três peças.
(b) Peça $5$ e vazio trocados: a posição é a
transposição dos conteúdos das casas $5$ e $9$, de modo que
$\varepsilon = -1$; o vazio fica no centro, à distância de Manhattan
$d = 2$ da sua casa, logo $(-1)^d = +1$ e $I = -1$:
inalcançável. Não é possível simplesmente ``estacionar o vazio no meio''
deixando as peças de resto ordenadas.

\textbf{23.} Seja $f \colon \mathfrak S_n \to \{\pm1\}$ um
\omterm{def:b1:structures:morphism}{morfismo}. Para duas transposições quaisquer $\tau, \tau'$, a questão 16
fornece $\sigma$ com $\sigma\tau\sigma^{-1} = \tau'$ (leve os
dois pontos movidos sobre os outros dois; $n \geq 3$ garante espaço
para fazê-lo, embora mesmo $n = 2$ seja trivial aqui). Então $f(\tau') =
f(\sigma)f(\tau)f(\sigma)^{-1} = f(\tau)$, pois $\{\pm1\}$ é
\omterm{def:b1:structures:group}{abeliano}: $f$ é constante nas transposições. Se essa constante é
$+1$, então $f = 1$ em todos os produtos de transposições, isto é,
em toda parte (questão 5). Se ela é $-1$, então $f(\sigma) =
(-1)^p = \varepsilon(\sigma)$ num produto de $p$ transposições.
Logo, $f \in \{1, \varepsilon\}$.

\textbf{24.} (i) A propriedade de \omterm{def:b1:structures:morphism}{morfismo} de $\varepsilon$ e a
maquinaria dos \omterm{def:b1:structures:morphism}{núcleos} deram a $\mathfrak A_n$ a sua estrutura de \omterm{def:b1:structures:subgroup}{subgrupo} e
o seu tamanho, e raciocínios ao estilo da \cref{prop:b1:structures:kernel}
percorrem as questões 11 e 21. (ii) Contagem:
$\abs{\mathfrak S_n} = n!$, o argumento de metade da questão 11
e a contagem $20\,160$ da questão 20 são o
\cref{ch:b1:counting} em ação. (iii) As questões 8--9 resolvem um
genuíno problema de boa definição --- ``a paridade do número
de transposições'' pressupõe que essa paridade não dependa
da fatoração, exatamente como as operações de $\Z/n\Z$ exigiram
independência do representante na \cref{def:b1:structures:zn}.

\textbf{25.} O sinal é um único cálculo com valores em $\{\pm1\}$,
demonstrado uma só vez ser bem definido, e ele faz três
trabalhos ao mesmo tempo: internamente, corta $\mathfrak S_n$ ao meio e
isola $\mathfrak A_n$ com os seus geradores $3$-ciclos;
externamente, decide numa linha uma questão (``estas duas
peças podem ser trocadas?'') que uma busca ingênua jamais resolveria,
pois nenhuma lista finita de sequências de movimentos fracassadas
demonstra impossibilidade; e, estruturalmente, é o motor de sinais
alternados dentro da fórmula $\det A = \sum_\sigma \varepsilon(\sigma)\,
a_{1\sigma(1)}\cdots a_{n\sigma(n)}$ do \cref{ch:b1:det}. O
padrão --- encontrar uma grandeza conservada por todo movimento
elementar, calculá-la no início e no alvo --- é a arma padrão do
matemático contra perguntas do tipo ``é possível?'',
e ele voltará sempre que um \omterm{def:b1:structures:group}{grupo} agir sobre um \omterm{def:b1:logic:sets}{conjunto} de
estados.
\end{solution}
